\chapter{Analisi del dominio e dei requisiti}

% ======================================================
%
% 1.1 - DESCRIZIONE DEL DOMINIO APPLICATIVO
%
% ======================================================
\section{Descrizione del dominio applicativo}
Il sistema analizzato modella il comportamento di un ascensore e dei suoi utilizzatori all'interno di un edificio a sviluppo verticale.

% ======================================================
% 1.1.1 - EDIFICIO E UTENTI
% ======================================================
\subsection{Edificio e utenti}
\label{subsec:building_users}
L'ambiente operativo è costituito da una struttura articolata su un numero finito $N$ di piani, identificati dall'insieme $\mathcal{F} = \{0, 1, \dots, N-1\}$.
All'interno di questa configurazione, il piano $0$ assume convenzionalmente il ruolo di piano terra (\emph{ground floor}).
Esso costituisce il punto cardine sia per l'accesso e l'uscita ordinaria dall'edificio, sia per le manovre di evacuazione rapida in situazioni di emergenza.

A popolare l'edificio è presente un insieme di utenti distribuiti lungo i vari livelli:
\begin{itemize}
    \item a ciascun piano possono trovarsi contemporaneamente fino a $M$ persone intenzionate a utilizzare l'impianto di risalita;
    \item ogni utente agisce come un agente autonomo e concorrente rispetto agli altri, caratterizzato dalla volontà di raggiungere un piano differente rispetto a quello in cui si trova attualmente.
\end{itemize}
Gli utenti dovranno quindi essere in grado di utilizzare l'ascensore per spostarsi all'interno dell'edificio.

% ======================================================
% 1.1.2 - CABINA E VINCOLO DI RISORSA
% ======================================================
\subsection{Cabina e vincolo di risorsa}
L'impianto è composto da una singola cabina addetta alla movimentazione verticale degli utenti tra i vari piani.
La caratteristica vincolante di tale risorsa risiede nella sua capacità strettamente unitaria ($Capacity = 1$): in nessun caso è consentito a più di una persona occupare lo spazio interno contemporaneamente.

L'interazione tra gli utenti e il sistema dell'ascensore avviene attraverso due distinte interfacce fisiche.
All'esterno della cabina, su ogni singolo piano, è presente un pulsante di chiamata (\emph{call button}) mediante il quale gli utenti in attesa segnalano la propria volontà di utilizzare l'ascensore.
All'interno della cabina è invece installata una pulsantiera dotata di tasti di selezione (\emph{floor buttons}), accessibile unicamente al passeggero salito a bordo per indicare il piano di destinazione verso cui essere trasportato.

% ======================================================
% 1.1.3 - POLITICA DI SERVIZIO FIFO
% ======================================================
\subsection{Politica di servizio FIFO}
Al fine di garantire equità di trattamento (\emph{fairness}) ed evitare situazioni in cui un utente venga ignorato o scavalcato indefinitamente (\emph{starvation}), gli utenti che richiedono l'ascensore su un determinato piano devono essere serviti rispettando l'ordine temporale di chiamata, attuando una politica di accodamento FIFO (\emph{First-In, First-Out}).

% ======================================================
%
% 1.2 - GESTIONE DELL'ALLARME ANTINCENDIO
%
% ======================================================
\section{Gestione dell'allarme antincendio}
\label{sec:allarme}
Il sistema dell'ascensore non opera completamente isolato, ma interagisce con un impianto antincendio esterno, in grado di generare in maniera asincrona un segnale di allarme in caso di incendio.

A fronte della ricezione del segnale di incendio, la tutela dell'utente prevale in modo assoluto sull'ordinario funzionamento dell'ascensore.
Il sistema attua pertanto la seguente sequenza di sicurezza:
\begin{enumerate}
    \item \textbf{Interruzione delle attività correnti}: l'ascensore interrompe le normali corse e sospende le chiamate registrate ai piani;
    \item \textbf{Discesa prioritaria al piano terra}: la cabina si dirige verso il piano terra ($Floor = 0$), ignorando qualsiasi fermata intermedia;
    \item \textbf{Apertura porte ed evacuazione}: una volta giunta al piano terra, le porte della cabina si aprono per permettere l'evacuazione dell'eventuale passeggero a bordo, garantendogli una via di fuga;
    \item \textbf{Stato di attesa cautelativa (\emph{Safe Lockout})}: l'ascensore rimane fermo al piano terra con le porte aperte, impedendo qualsiasi comando di azionamento;
    \item \textbf{Ripristino del servizio}: il blocco rimane attivo fino alla ricezione di un segnale esterno di risoluzione dell'allarme (\emph{alarm resolved}), a seguito del quale il sistema è autorizzato a ripristinare l'operatività.
\end{enumerate}

% ======================================================
%
% 1.3 - REQUISITI E PROPRIETÀ DI CORRETTEZZA FORMALE
%
% ======================================================
\section{Requisiti e proprietà di correttezza formale}
Sulla base della specifica del problema, il sistema concorrente deve soddisfare le proprietà di \emph{Safety} e \emph{Liveness}.

% ======================================================
% 1.3.1 - PROPRIETÀ DI SAFETY (SICUREZZA)
% ======================================================
\subsection{Proprietà di Safety (Sicurezza)}
Le proprietà di sicurezza sanciscono formalmente che durante l'intera evoluzione temporale del sistema ``non deve mai verificarsi alcuno stato non consentito''.

Nel contesto dell'ascensore, il requisito primario di sicurezza coincide con la \textbf{mutua esclusione della cabina} (\emph{mutual exclusion}).
In nessun istante temporale e lungo alcuna traiettoria di esecuzione possono trovarsi due o più persone simultaneamente all'interno dell'ascensore.
Indicando con $\text{Utenti\_in\_cabina}(t)$ il numero di occupanti a bordo all'istante $t$, il controllore deve garantire costantemente che:
\[
  \forall t \ge 0, \quad \text{Utenti\_in\_cabina}(t) \le 1
\]

% ======================================================
% 1.3.2 - PROPRIETÀ DI LIVENESS (VIVACITÀ CONDIZIONATA)
% ======================================================
\subsection{Proprietà di Liveness (Vivacità condizionata)}
Mentre le proprietà di sicurezza vietano configurazioni illegali, le proprietà di liveness stabiliscono che gli eventi desiderati dal sistema ``debbano prima o poi verificarsi nel tempo''.
Nel nostro dominio tali proprietà possono essere violate in un particolare caso, in quanto l'arrivo di un allarme antincendio mette indefinitamente in pausa le promesse di servizio ordinario a favore della sicurezza degli utenti (fino alla cessazione dell'allarme):
\begin{itemize}
    \item \textbf{Arrivo al piano di chiamata (\emph{Eventual Arrival upon Call})}: a fronte della pressione del pulsante di chiamata da parte di un utente a un dato piano, la cabina deve prima o poi giungere a quel piano per consentire il servizio, a meno che non si manifesti uno stato di allarme antincendio;
    \item \textbf{Consegna al piano desiderato (\emph{Eventual Delivery})}: quando un passeggero a bordo seleziona il piano prescelto, l'ascensore deve raggiungere tale destinazione, salvo il caso in cui un segnale di allarme costringa la cabina a dirigersi verso il piano terra per l'evacuazione dell'utente;
    \item \textbf{Accesso al servizio ed assenza di starvation (\emph{Starvation Freedom})}: ciascun utente posizionato in attesa a un piano deve infine riuscire a salire a bordo della cabina, escludendo il verificarsi di un'emergenza antincendio che impedisce l'utilizzo dell'ascensore.
\end{itemize}